\section{Conclusion}
\label{sec:conclusion}

This paper has specified, in full, the methodology of \tool{} at methodology
version \methver{}: every input and its normalisation, the fourteen-typology
evidence library and the derivations behind it, every scoring formula and every
weight, the treatment of uncertainty and missing data, the design of the
sensitivity analysis, and the engineering arrangement that keeps the software
in correspondence with the specification.

The instrument occupies a deliberately narrow position. It is faster and more
comparable than unstructured judgment, and far less capable than microclimate
simulation. It exists for the triage step: narrowing a long list of candidate
sites and interventions to a short one, on a transparent and
evidence-referenced basis, and recording why. Judged as a simulation it fails;
judged as a screening instrument, the questions that matter are whether its
values are traceable, whether its uncertainties are honestly represented,
whether its rankings are stable, whether its value judgments are disclosed, and
whether it declines to answer when its inputs do not support an answer
(\Cref{sec:standards}).

Several of the methodology's choices lower its apparent authority, and they
were the substance of the work rather than concessions to it. Adopted cooling
envelopes sit at or below synthesis central estimates because much of the
underlying evidence is modelling-derived and modelling reports larger effects
than measurement. The mixed intervention package is capped at the
best-evidenced single-measure ceiling because no retrieved source quantifies
super-additive cooling. Reported temperature ranges are clipped to the
literature envelope, so that a favourable site raises a relative score but never
a physical claim. Per-typology energy factors that no source supported were
deleted and replaced by a derivation from a published sensitivity. No default
unit costs ship, so the tool declines to produce the payback figure that would
otherwise be its most quotable output. Green fa\c{c}ades, rain gardens, and
courtyard greening are labelled low confidence and their confidence caps
propagate into results.

The limitations are correspondingly explicit. All values are daytime, and
nighttime heat --- the more consequential quantity for health outcomes --- is
not estimated. The quantity modelled is air temperature, not thermal comfort,
which means the tool systematically understates shade-dominant interventions;
the bias is stated rather than corrected by an unsourced factor. The evidence
base is skewed toward temperate and subtropical Northern Hemisphere cities, and
the confidence model does not yet measure that skew. The aggregation weights are
expert judgment, including a disclosed equity-forward stance that gives
vulnerability an effective weight of \num{0.25} against heat exposure at
\num{0.15}. The sensitivity analysis of \Cref{sec:sensitivity-results} finds the ranking
substantially stable under \SI{\pm 25}{\percent} single-weight perturbation ---
worst-case pairwise stability \num{0.9737}, with three near-boundary category
changes in 240 scenario-perturbation combinations --- which bounds how much the
contestable weights actually move decisions.

The methodology was first offered for review with its gaps disclosed rather
than closed, because the values most in need of scrutiny are settled most
cheaply before an engine, a user base, and a set of published assessments have
accumulated around them. Version \methver{} closes the disclosed specification
gaps --- the sub-indicator rules, the cost-feasibility derivation, and the
confidence field mapping --- without altering any evidence-derived value. The tool, its complete configuration, and this
document are public under the Apache License 2.0. Every value it applies names
its source; \Cref{app:verification} states what was actually verified about
each of those sources, including where publisher restrictions meant a finding
was confirmed from an abstract or an indexing record rather than a full text.
That disclosure is not a weakness in the evidence base. It is the property that
makes the evidence base checkable, and it is what the authors would want to
know before trusting an instrument of this kind.
